\section{RESULTADOS}

El análisis molecular abarcó tres conjuntos de muestras de la PTAR Quitumbe. Para SARS-CoV-2 se procesaron 42 muestras semanales recolectadas entre las semanas 1 y 42 (14 de abril de 2025 al 1 de febrero de 2026). Para enterovirus y para RSV se analizaron 22 muestras cada uno, conformadas por 14 muestras prospectivas (semanas 43 a 56) y 8 muestras retrospectivas conservadas (semanas 17 a 20 y 34 a 37). El rendimiento de la secuenciación fue desigual entre los tres virus y entre los tipos de muestra (Tabla~\ref{tab:resumen_rendimiento}, Figura~\ref{fig:f3_rendimiento}).

De las 42 muestras de SARS-CoV-2, 21 se sometieron a secuenciación y 6 permitieron asignar clado y linaje; en las 15 restantes la biblioteca aportó entre 0 y 1 lectura de SARS-CoV-2, de modo que el consenso no cubrió ninguna posición del genoma (Anexos~\ref{anexo:a} y~\ref{anexo:b}). En enterovirus, 2 de las 22 muestras produjeron secuencia de la región VP1, tipificada como Coxsackievirus A1, y ninguna rindió secuencia de poliovirus. En RSV, 4 muestras se secuenciaron y ninguna aportó lecturas del virus. Las muestras retrospectivas conservadas, analizadas para enterovirus y para RSV, no rindieron secuencia en ninguno de los dos métodos de preservación empleados (Figura~\ref{fig:f4_conservacion}).

\begin{longtblr}[
  caption={Resumen del esfuerzo y el rendimiento de la secuenciación por virus},
  label={tab:resumen_rendimiento}
]{
  width=\textwidth,
  colspec={X[1.3,l] X[1,c] X[1.1,c] X[1.2,c] X[1.2,c] X[1.2,c]},
}
\textbf{Virus} & \textbf{Muestras} & \textbf{Secuenc.} & \textbf{Con caracterización} & \textbf{Sin caracterización} & \textbf{No secuenc.} \\
SARS-CoV-2 & 42 & 21 & 6 & 15 & 21 \\
Enterovirus & 22 & 2 & 2\textsuperscript{*} & 0 & 20 \\
RSV & 22 & 4 & 0 & 4 & 18 \\
\end{longtblr}

\noindent\footnotesize{\textsuperscript{*}Las dos muestras caracterizadas correspondieron a Coxsackievirus A1, un enterovirus no polio de la especie C tipificado por la secuencia de VP1; no se obtuvo secuencia de poliovirus.}\normalsize

\begin{figure}[H]
  \centering
  \includegraphics[width=0.62\textwidth]{resultados/f3_rendimiento.pdf}
  \caption{Esfuerzo y rendimiento de la secuenciación por virus. Las barras muestran el número de muestras no secuenciadas, secuenciadas sin caracterización y secuenciadas con caracterización (Coxsackievirus A1 en el caso de \textit{enterovirus}).}
  \label{fig:f3_rendimiento}
\end{figure}

\subsection{Detección molecular por RT-qPCR}

La RT-qPCR detectó el gen N1 de SARS-CoV-2 en 31 de las 42 muestras semanales (73,8\,\%; IC 95\,\%: 58,9--84,7\,\%, método de Wilson) y el gen N2 en 21 (50,0\,\%; IC 95\,\%: 35,5--64,5\,\%), con valores de Cq tardíos propios de aguas residuales con baja carga viral y ARN fragmentado (Figura~\ref{fig:f5_cq}) \cite{Parkins_2023,Child_2023}. Las once muestras sin Cq de N1 corresponden a las semanas 1 y 2, que solo pasaron por el primer análisis, en el que ninguna muestra amplificó; a ocho semanas sin amplificación de N1 ni de N2, y a la semana 33, en la que amplificó únicamente N2 (Anexo~\ref{anexo:a}). El Cq de N1 difirió entre los tres resultados de la secuenciación (Kruskal-Wallis, $H=16{,}29$; gl $=2$; $p<0{,}001$) (Tabla~\ref{tab:cq_resumen}). Las seis muestras con linaje asignado presentaron la mediana más baja (34,0; RIC 33,4--34,1), seguidas por las 14 muestras secuenciadas sin caracterización (36,4; RIC 36,1--37,1) y por las 11 muestras no secuenciadas con Cq registrado (38,5; RIC 37,1--39,1), y las tres comparaciones por pares conservaron la significación tras la corrección de Bonferroni (Mann-Whitney, $p$ ajustado entre 0,001 y 0,033). Los rangos de los tres grupos se solapan, ya que las semanas 13 y 14, con Cq de N1 de 33,9 y 34,9, no aportaron lecturas de SARS-CoV-2, mientras que la semana 3, con 35,8, permitió asignar linaje; por ello el Cq acompaña la probabilidad de recuperar información genómica sin determinarla por sí solo.

\begin{longtblr}[
  caption={Valores de Cq del gen N1 de SARS-CoV-2 según el resultado de la secuenciación},
  label={tab:cq_resumen}
]{
  width=\textwidth,
  colspec={X[2,l] X[0.6,c] X[1.3,c] X[1,c]},
}
\textbf{Resultado de la secuenciación} & \textbf{n} & \textbf{Cq N1, mediana (RIC)} & \textbf{Cq N1, rango} \\
Linaje determinado & 6 & 34,0 (33,4--34,1) & 33,1--35,8 \\
Secuenciado, sin caracterización & 14 & 36,4 (36,1--37,1) & 33,9--42,3 \\
No secuenciado (con Cq registrado) & 11 & 38,5 (37,1--39,1) & 34,4--42,9 \\
\textbf{Total con Cq de N1} & \textbf{31} & & \\
\end{longtblr}

\noindent\footnotesize{De las 42 muestras semanales, 31 aportaron Cq de N1. Las 11 restantes corresponden a las semanas 1 y 2 (análisis 1, sin amplificación de ninguna muestra), a las semanas 15, 18, 19, 20, 22, 26, 27 y 31 (sin amplificación de N1 ni de N2) y a la semana 33 (solo N2). La semana 31 se secuenció sin Cq de N1 ni de N2.}\normalsize

\begin{figure}[H]
  \centering
  \includegraphics[width=\textwidth]{resultados/f5_cq_sars.pdf}
  \caption{Valores de Cq de SARS-CoV-2 por semana para los genes N1 (HEX) y N2 (FAM). El eje vertical está invertido, de modo que los valores de Cq más bajos, que corresponden a mayor carga viral, aparecen en la parte superior. El color indica el resultado de la secuenciación; la banda verde señala el rango de Cq de las muestras con linaje determinado.}
  \label{fig:f5_cq}
\end{figure}

\subsection{Controles de laboratorio}

El control positivo del ensayo Luna amplificó en los ocho análisis de RT-qPCR, con Cq de N1 entre 20,38 y 25,52 y de N2 entre 20,03 y 24,47, y el control sin molde (NTC) no presentó señal en N1 ni en N2 en ninguno de ellos (Anexo~\ref{anexo:c}). Los hisopados faríngeos empleados como control positivo de extracción en los análisis 2 a 7 rindieron Cq de N1 entre 11,28 y 20,83, lo que verifica la extracción y la transcripción reversa del ARN de SARS-CoV-2; el primer análisis se ejecutó sin ese control y ninguna muestra amplificó, por lo que sus resultados se descartaron. En el análisis 8, el NTC presentó una señal tardía del control interno RP (Cq 40,84), sin señal de N1 ni de N2. El control interno RP, que detecta ARN humano, amplificó en 19 de las 40 muestras de los análisis 2 a 8, según el análisis conservado para cada semana.

Los controles de la PCR multiplex se evaluaron en gel y no se secuenciaron, de modo que ninguna corrida de secuenciación incluyó un control negativo con código de barras. Sin ese control no es posible estimar la contaminación cruzada entre muestras ni la asignación errónea de lecturas entre códigos de barras.

\subsection{Linajes de SARS-CoV-2}

La caracterización genómica de SARS-CoV-2 se limitó a las seis muestras de las semanas 3 a 9 (28 de abril al 15 de junio de 2025), todas procesadas en la primera corrida de secuenciación. Los consensos generados por EPI2ME cubrieron entre el 3,1 y el 17,2\,\% del genoma de referencia y mostraron entre 3 y 14 sustituciones respecto de él, de modo que Nextclade les otorgó una calificación de calidad deficiente y Pangolin no les asignó linaje (Tabla~\ref{tab:linajes_sars}). Con esa información parcial, Nextclade ubicó los consensos de las semanas 3 a 8 en el clado 24A, con los linajes JN.1, LZ.5 y LU.2.1.1, y el de la semana 9 en el clado 24H, con el linaje LF.7.1.3; todos ellos descienden de JN.1.

La asignación basada en el consenso resultó inestable entre réplicas y versiones del análisis. Las réplicas de las semanas 7, 8 y 9 amplificadas con VarSkip recibieron el linaje JN.1, distinto del asignado a las amplificadas con ARTIC, y los mismos consensos analizados con la versión de septiembre de 2026 del conjunto de datos de Nextclade cambiaron de linaje en cuatro de las seis semanas.

\begin{longtblr}[
  caption={Clados y linajes de SARS-CoV-2 asignados con Nextclade a los consensos de EPI2ME},
  label={tab:linajes_sars}
]{
  width=\textwidth,
  colspec={X[0.8,c] X[1.2,c] X[0.8,c] X[1.1,c] X[1.3,c] X[1.3,c]},
}
\textbf{Semana} & \textbf{Fecha} & \textbf{Clado} & \textbf{Linaje} & \textbf{Cobertura del consenso (\%)} & \textbf{Linaje de la réplica VarSkip} \\
3 & 28/04/2025 & 24A & JN.1 & 3,1 & -- \\
5 & 12/05/2025 & 24A & LZ.5 & 4,5 & -- \\
6 & 19/05/2025 & 24A & JN.1 & 4,9 & JN.1 \\
7 & 26/05/2025 & 24A & LZ.5 & 4,6 & JN.1 \\
8 & 02/06/2025 & 24A & LU.2.1.1 & 5,8 & JN.1 \\
9 & 09/06/2025 & 24H & LF.7.1.3 & 17,2 & JN.1 \\
\end{longtblr}

\noindent\footnotesize{Asignación con el conjunto de datos de Nextclade del 9 de junio de 2025 sobre los consensos amplificados con ARTIC V3. La cobertura corresponde a la fracción del genoma con bases determinadas según Nextclade. Todos los consensos obtuvieron calificación de calidad deficiente en Nextclade y Pangolin no les asignó linaje.}\normalsize

La deconvolución con Freyja, sobre una fracción del genoma cubierta a 10 lecturas, tras recortar los cebadores, de entre el 20,8 y el 42,1\,\%, atribuyó entre el 86 y el 97\,\% de la mezcla de las diez bibliotecas a QA.4, un sublinaje de LF.7.1.3 del clado 24H, y hasta el 13\,\% a LZ.5 en siete de ellas (Tabla~\ref{tab:freyja}, Figura~\ref{fig:f2_linajes}). Los linajes minoritarios detectados por Freyja coinciden con los que el consenso asignó en las semanas 5, 7 y 8 (LZ.5 y un linaje afín a LU.2). La composición se mantuvo semejante entre semanas, entre esquemas de cebadores y con la biblioteca de mutaciones actual, de modo que los datos describen la presencia conjunta de descendientes de JN.1 durante el periodo caracterizado y no sostienen una sucesión temporal de linajes. Esa uniformidad entre bibliotecas procesadas en una misma corrida, sin un control negativo secuenciado, tampoco permite descartar la contaminación cruzada entre muestras.

\begin{longtblr}[
  caption={Abundancia relativa de linajes de SARS-CoV-2 estimada con Freyja por biblioteca},
  label={tab:freyja}
]{
  width=\textwidth,
  colspec={X[0.7,c] X[1,c] X[1.1,c] X[0.8,c] X[0.8,c] X[0.8,c] X[1.2,c]},
}
\textbf{Semana} & \textbf{Cebadores} & \textbf{Genoma $\geq$10x (\%)} & \textbf{QA.4 (\%)} & \textbf{LZ.5 (\%)} & \textbf{Otros (\%)} & \textbf{Linaje del consenso} \\
3 & ARTIC V3 & 23,9 & 94 & 6 & 0 & JN.1 \\
5 & ARTIC V3 & 28,6 & 93 & 6 & 1 & LZ.5 \\
6 & ARTIC V3 & 27,5 & 94 & 3 & 2 & JN.1 \\
6 & VarSkip & 20,8 & 86 & 13 & 1 & JN.1 \\
7 & ARTIC V3 & 33,0 & 94 & 5 & 1 & LZ.5 \\
7 & VarSkip & 21,7 & 97 & 3 & 0 & JN.1 \\
8 & ARTIC V3 & 31,8 & 95 & 0 & 5 & LU.2.1.1 \\
8 & VarSkip & 24,4 & 91 & 7 & 1 & JN.1 \\
9 & ARTIC V3 & 42,1 & 95 & 0 & 5 & LF.7.1.3 \\
9 & VarSkip & 24,2 & 89 & 0 & 11 & JN.1 \\
\end{longtblr}

\noindent\footnotesize{Freyja 2.0.5 con la biblioteca de mutaciones definitorias del 20 de junio de 2025 y profundidad mínima de 10 lecturas. «Otros» agrupa linajes con abundancia individual de hasta el 6\,\%; en la semana 8 (ARTIC V3) corresponde a un linaje afín a LU.2. Los porcentajes se redondearon a la unidad. La cobertura se calculó sobre las lecturas con los cebadores recortados, por lo que difiere ligeramente de la del Anexo~\ref{anexo:b}.}\normalsize

\begin{figure}[H]
  \centering
  \includegraphics[width=\textwidth]{resultados/f2_linajes_sars.pdf}
  \caption{Abundancia relativa de linajes de SARS-CoV-2 estimada con Freyja en las diez bibliotecas de las semanas 3 a 9. Sobre cada barra se indica el linaje asignado al consenso con Nextclade.}
  \label{fig:f2_linajes}
\end{figure}

Las semanas con linaje determinado se ubicaron entre los valores de Cq más bajos del periodo (Figura~\ref{fig:f5_cq}). A partir de la semana 10 se secuenciaron 15 muestras más en tres corridas, y ninguna aportó cobertura del genoma de SARS-CoV-2, ya que de las 11\,387 a 173\,030 lecturas obtenidas por biblioteca solo entre 0 y 1 correspondieron al virus (Anexo~\ref{anexo:b}). En la segunda corrida, entre el 28,9 y el 75,3\,\% de las lecturas portaba en sus extremos cebadores del esquema de SARS-CoV-2 sin alinear con el genoma viral, lo que identifica productos de amplificación inespecífica sobre ácidos nucleicos ajenos al virus; en las dos corridas siguientes esa proporción bajó a 2,2--19,5\,\%, y la biblioteca quedó compuesta casi por completo por material no amplificado por los cebadores. El Cq de N1 se correlacionó con la fracción del genoma cubierta a 10 lecturas en las 20 muestras secuenciadas con Cq registrado (Spearman, $\rho=-0{,}70$; $p<0{,}001$), aunque esa relación coincide con el orden de las corridas, pues todas las bibliotecas sin lecturas del virus proceden de corridas posteriores a la semana 9.

\subsection{\textit{Enterovirus}}

El flujo de PCR anidada recuperó material de enterovirus en las muestras de las semanas 45 y 46 (febrero de 2026), cuyos productos de la segunda ronda mostraron bandas de 1100 a 1200 pares de bases y se secuenciaron en duplicado (Tabla~\ref{tab:polio_rsv_conservacion}). En el análisis inicial, las lecturas se alinearon contra referencias de poliovirus de los tipos 1, 2 y 3 y solo se ubicaron entre las posiciones 529 y 1049, en la región VP4, lo que llevó a interpretar que los cebadores de VP1 habían amplificado otra región. El reanálisis contra un panel de enterovirus de la especie C corrigió esa interpretación, ya que el 91,5\,\% de las lecturas de la semana 45 y el 51,7\,\% de las de la semana 46 alineadas con los genomas completos de referencia correspondieron a la región VP1, con más del 97\,\% del amplicón de 1094 pb cubierto a 20 lecturas en las cuatro bibliotecas (Anexo~\ref{anexo:b}). Ninguna de esas lecturas alineó con las referencias de poliovirus, que solo recuperaron las de la región 5'NTR--VP4, más conservada, y por ello el análisis inicial registró únicamente esa región.

Las secuencias consenso de VP1 de las dos semanas resultaron idénticas en 1087 pb, y su mayor identidad en el GenBank fue del 87,1\,\% con un Coxsackievirus A1 aislado en 2023 (PQ566821.1; valor E de 0); las 20 secuencias de mayor identidad correspondieron también a Coxsackievirus A1. Esa identidad supera el 75\,\% que debe alcanzar la secuencia prototipo o, en enterovirus de la especie C, el aislado del GenBank de mayor identidad para asignar el tipo por la secuencia de VP1 \cite{Martin_2016}, y el árbol de máxima verosimilitud ubicó ambas secuencias dentro del clado de Coxsackievirus A1, con un apoyo del 100\,\% en las réplicas de bootstrap ultrarrápido (Figura~\ref{fig:arbol_vp1}). Las lecturas de la región 5'NTR--VP4--VP2 coincidieron con distintos tipos de la especie C, incluidos algunos poliovirus, pero el tipo se establece por la secuencia de VP1 \cite{Martin_2016}. La identidad completa entre las secuencias de las dos semanas es compatible con la circulación de una misma cepa en semanas consecutivas, aunque ambas muestras se amplificaron y secuenciaron juntas, sin un control negativo secuenciado que descarte la contaminación cruzada.

\begin{figure}[H]
  \centering
  \includegraphics[width=0.85\textwidth]{resultados/f6_arbol_vp1.pdf}
  \caption{Árbol de máxima verosimilitud de la región VP1 (1094 pb) de enterovirus de la especie C, con enterovirus A como grupo externo. En rojo, los consensos de las semanas 45 y 46. Los números indican el apoyo de bootstrap ultrarrápido ($\geq$90). Las ramas marcadas con // se acortaron en el dibujo.}
  \label{fig:arbol_vp1}
\end{figure}

Dentro de este género, ninguna de las 22 muestras rindió secuencia de poliovirus.

\subsection{Virus sincitial respiratorio}

El RSV se intentó amplificar en las 22 muestras y ninguna rindió secuencia del virus. Las muestras de las semanas 50, 52, 53 y 54 (marzo y abril de 2026) mostraron bandas en la electroforesis y se secuenciaron en ocho bibliotecas, una con ARN sin diluir y otra con ARN diluido por muestra, que reunieron 581\,826 lecturas (Anexo~\ref{anexo:b}). Ninguna de esas lecturas alineó con calidad de mapeo de al menos 20 con las referencias de RSV-A o RSV-B, en concordancia con el análisis inicial con EPI2ME, que no encontró alineamiento con la referencia de RSV (Tabla~\ref{tab:polio_rsv_conservacion}). Entre el 73,2 y el 78,4\,\% de las lecturas portaba en sus extremos cebadores del esquema ARTIC RSV, con una longitud mediana de 155 a 234 pb, inferior a los cerca de 400 pb de los amplicones del esquema \cite{Maloney_2025}, de modo que las bandas que motivaron la secuenciación correspondieron a productos de amplificación inespecífica. Como el RSV no se analizó por RT-qPCR, el estudio no dispone de una medida de positividad ni de carga viral para este virus, y el resultado constituye un intento de detección negativo.

\begin{longtblr}[
  caption={Muestras de enterovirus y RSV por tipo de muestra, método de conservación y desenlace},
  label={tab:polio_rsv_conservacion}
]{
  width=\textwidth,
  colspec={X[1,l] X[1.2,l] X[1.3,l] X[0.7,c] X[0.9,c] X[1.5,l]},
}
\textbf{Virus} & \textbf{Tipo} & \textbf{Conservación} & \textbf{n} & \textbf{Secuenc.} & \textbf{Desenlace} \\
\SetCell[r=3]{m} Enterovirus & Prospectiva & $-20$\,°C & 14 & 2 & Coxsackievirus A1 (sem. 45 y 46) \\
Enterovirus & Retrospectiva & Shield 2X & 4 & 0 & Sin secuencia \\
Enterovirus & Retrospectiva & PBS 1X & 4 & 0 & Sin secuencia \\
\SetCell[r=3]{m} RSV & Prospectiva & $-20$\,°C & 14 & 4 & Sin lecturas de RSV (sem. 50, 52--54) \\
RSV & Retrospectiva & Shield 2X & 4 & 0 & Sin secuencia \\
RSV & Retrospectiva & PBS 1X & 4 & 0 & Sin secuencia \\
\end{longtblr}

\begin{figure}[H]
  \centering
  \includegraphics[width=\textwidth]{resultados/f4_conservacion.pdf}
  \caption{Muestras de enterovirus y RSV según el método de conservación. Las muestras prospectivas, congeladas a $-20$\,°C hasta su procesamiento, concentraron la totalidad de las secuenciaciones; las muestras retrospectivas conservadas en Shield 2X y en PBS 1X no rindieron secuencia.}
  \label{fig:f4_conservacion}
\end{figure}

\subsection{Comparación temporal entre virus}

La distribución temporal de los resultados muestra ventanas de análisis con escaso solapamiento entre los tres virus (Figura~\ref{fig:f1_heatmap}). La caracterización de SARS-CoV-2 se concentró entre las semanas 3 y 9 de 2025, seguida de un periodo de secuenciación sin cobertura del genoma hasta la semana 32. SARS-CoV-2 no se analizó entre las semanas 43 y 56, periodo en el que el flujo de trabajo se orientó a los enterovirus y al RSV. El análisis de enterovirus y RSV se ubicó en el tramo prospectivo de 2026 y en las semanas retrospectivas de agosto y diciembre de 2025. Los linajes de SARS-CoV-2 se obtuvieron en 2025 y Coxsackievirus A1 en febrero de 2026; las bibliotecas de RSV, de marzo y abril de 2026, no contuvieron lecturas del virus. 

\begin{figure}[H]
  \centering
  \includegraphics[width=\textwidth]{resultados/f1_heatmap_temporal.pdf}
  \caption{Estado de detección y secuenciación de SARS-CoV-2, enterovirus y RSV por semana de muestreo. El panel superior abarca las semanas 1 a 28 y el inferior las semanas 29 a 56. Cada celda indica si el virus no se analizó, no se secuenció, se secuenció sin caracterización o se caracterizó (linaje en SARS-CoV-2, Coxsackievirus A1 en enterovirus).}
  \label{fig:f1_heatmap}
\end{figure}
